\chapter{FIO Library}

The FIO library is the assembly programmer's bridge to NovaHost storage and
host-backed services. It is not a POSIX-style file API and it is not a normal
6502 disk controller. It is a small command mailbox at \novacode{\$B9A0-\$B9EF}
that NovaHost watches, executes, and reports through status registers.

\begin{novanote}
Use \novacode{fio.inc} for declarations and argument aliases. Include
\novacode{fio.s} after the code that calls it. The implementation emits only
referenced routines by default, so including it before the call sites can strip
the routines you expected to link.
\end{novanote}

\begin{lstlisting}
.include "fio.inc"

      ; program code that calls fio_load, fio_gload, etc.

.include "fio.s"
\end{lstlisting}

Define \novacode{NOVA\_EMIT\_ALL\_RUNTIME} or
\novacode{FIO\_EMIT\_ALL\_RUNTIME} before including \novacode{fio.s} when a ROM
or debugging build should keep the whole implementation.

\section{What FIO Is}

FIO is a hardware-facing command protocol plus a small set of assembly helpers.
Your program fills filename, address, length, and mode registers, then issues a
command. NovaHost performs the storage operation and writes back status, error,
size, and directory metadata.

\begin{center}
\novasvg[width=\textwidth]{build/diagrams/fio-architecture}
\end{center}

There are four pieces to keep straight:

\begin{description}
\item[Command register] \novacode{FIO\_CMD} triggers work when written. The
library helpers use the command constants from \novacode{nova.inc}.
\item[Status and error registers] \novacode{FIO\_STATUS} becomes
\novacode{FIO\_STATUS\_OK} or \novacode{FIO\_STATUS\_ERROR}. On error,
\novacode{FIO\_ERRCODE} distinguishes not found, I/O failure, end of directory,
disk full, and not mounted.
\item[Filename buffer] \novacode{FIO\_NAME} holds up to 63 bytes and
\novacode{FIO\_NAMELEN} tells the host how many bytes are valid.
\item[Argument registers] The same register bytes are reused by different
commands. For example, \novacode{FIO\_SRCL/H} can be a CPU address for
\novacode{fio\_load}, RNG result bytes after \novacode{fio\_rng}, or the low 16
bits of an \novacode{XPAGE} file offset.
\end{description}

This reuse is intentional. Treat FIO registers as a command mailbox, not as
persistent application state.

\section{Command Lifecycle}

Most callers should use \novacode{fio\_exec} or one of the named wrappers such
as \novacode{fio\_load}. \novacode{fio\_exec} clears the previous status and
error, writes the command byte, waits until the host reports completion, and
returns the shared runtime result: A is 0 for success and 1 for error.

\begin{center}
\novasvg[width=\textwidth]{build/diagrams/fio-command-flow}
\end{center}

\begin{center}
\small
\begin{tabularx}{\textwidth}{@{}l>{\raggedright\arraybackslash}X@{}}
\toprule
\textbf{Helper} & \textbf{Use} \\
\midrule
\novacode{fio\_issue} & Write A to \novacode{FIO\_CMD} and return without
waiting or checking status. Use this only when the caller owns the polling. \\
\novacode{fio\_exec} & Issue A as a command, wait for completion, and return
A=0 or A=1. \\
\novacode{fio\_check} & Convert the current \novacode{FIO\_STATUS} value into
the shared A=0/A=1 result. \\
\novacode{fio\_clear\_error} & Clear the host-visible error latch before a retry
or before returning to UI code that displays status. \\
\bottomrule
\end{tabularx}
\end{center}

The named wrappers all jump through \novacode{fio\_exec}. A failed command
returns A=1, but the useful reason is still in \novacode{FIO\_ERRCODE}. Check it
before preparing another FIO command, because the next command reuses the same
registers.

\section{Filename Handling}

Every file command reads \novacode{FIO\_NAME} and \novacode{FIO\_NAMELEN}.
Programs can write those registers directly, but pointer-based code should use
\novacode{fio\_copy\_name}. It validates the length range 1..63, copies bytes
from \novacode{FIO\_ARG\_NAMEPTR\_L/H}, and returns A=0 on success.

\begin{lstlisting}
copy_level_name:
      LDA   #<level_name
      STA   FIO_ARG_NAMEPTR_L
      LDA   #>level_name
      STA   FIO_ARG_NAMEPTR_H
      LDA   #level_name_end - level_name
      STA   FIO_ARG_NAMELEN
      JSR   fio_copy_name
      RTS

level_name:
      .byte "LEVEL1"
level_name_end:
\end{lstlisting}

The host accepts simple host filenames and mounted-device paths. Explicit
extensions matter for data files: if no extension is supplied, RAM program
loads look for \novacode{.bas} and then \novacode{.bin}; graphics transfers use
\novacode{.gfx}; XRAM transfers use \novacode{.xram}; paged slices use
\novacode{.bin}.

\section{CPU RAM Loads And Saves}

\novacode{fio\_load} moves a \novacode{.bas} or \novacode{.bin} file into CPU
RAM. For \novacode{.bas}, the destination comes from \novacode{FIO\_SRCL/H}.
For \novacode{.bin}, the file's two-byte load-address prefix supplies the
destination and NovaHost writes that address back to \novacode{FIO\_SRCL/H}.
The transferred payload size is returned in \novacode{FIO\_SIZEL/H}.

\begin{lstlisting}
load_level:
      LDA   #<level_name
      STA   FIO_ARG_NAMEPTR_L
      LDA   #>level_name
      STA   FIO_ARG_NAMEPTR_H
      LDA   #level_name_end - level_name
      STA   FIO_ARG_NAMELEN
      JSR   fio_copy_name
      BNE   @error

      LDA   #<$6000
      STA   FIO_SRCL
      LDA   #>$6000
      STA   FIO_SRCH
      JSR   fio_load
      BNE   @error
      RTS

@error:
      LDA   FIO_ERRCODE
      RTS

level_name:
      .byte "LEVEL1"
level_name_end:
\end{lstlisting}

\novacode{fio\_save} uses \novacode{FIO\_SRCL/H} as the first CPU address and
\novacode{FIO\_ENDL/H} as the exclusive end address. NovaHost writes a two-byte
load-address prefix before the saved bytes. A default \novacode{.bas} save still
loads to the caller-supplied \novacode{FIO\_SRCL/H}; an explicitly saved
\novacode{.bin} file uses the prefix as its load address.

\section{VGC, XRAM, And Paged Transfers}

Real Nova applications should avoid byte-copying large assets through the CPU
when a bulk path exists. FIO can load directly into VGC memory, stream whole
files to XRAM, or page slices of a file into XRAM, CPU RAM, or VGC memory.

\begin{center}
\novasvg[width=\textwidth]{build/diagrams/fio-transfer-paths}
\end{center}

\begin{center}
\small
\begin{tabularx}{\textwidth}{@{}l>{\raggedright\arraybackslash}X@{}}
\toprule
\textbf{Path} & \textbf{When to use it} \\
\midrule
\novacode{fio\_gload}, \novacode{fio\_gsave} & Bulk load or save VGC memory.
Set \novacode{FIO\_GSPACE}, \novacode{FIO\_GADDRL/H}, and
\novacode{FIO\_GLENL/H}. A zero load length means the whole file. \\
\novacode{FIO\_CMD\_XLOAD}, \novacode{FIO\_CMD\_XSAVE} & Stream a file to or
from flat XRAM. Use \novacode{fio\_prepare\_xram\_transfer} to copy the
pointer filename, 24-bit XRAM address, and length into FIO registers, then call
\novacode{fio\_run}. XLOAD accepts length zero for the whole file; XSAVE
requires a nonzero length. \\
\novacode{FIO\_CMD\_XPAGE} & Load a slice of a \novacode{.bin} file. The pager
library wraps this as \novacode{pager\_load\_file\_page}; target 0 is XRAM,
target 1 is CPU RAM, and target 2 is VGC memory. \\
\bottomrule
\end{tabularx}
\end{center}

For VGC transfers, \novacode{FIO\_GSPACE} is the VGC memory space and
\novacode{FIO\_GADDRL/H} is the 16-bit offset inside that space. For XRAM
transfers, the same bytes become a 24-bit address:
\novacode{FIO\_GADDRL}, \novacode{FIO\_GADDRH}, then
\novacode{FIO\_GSPACE}.

\section{Directory And Filesystem Commands}

\novacode{fio\_dir\_open} opens a listing and populates the first entry.
\novacode{fio\_dir\_read} advances to the next entry. When there are no more
entries, the host reports \novacode{FIO\_STATUS\_ERROR} with
\novacode{FIO\_ERR\_EOD}; that is normal iteration termination, not a broken
device.

\begin{center}
\small
\begin{tabularx}{\textwidth}{@{}l>{\raggedright\arraybackslash}X@{}}
\toprule
\textbf{Result register} & \textbf{Directory meaning} \\
\midrule
\novacode{FIO\_DIRTYPE} & Entry type: 0 for BASIC/default, 1 for SID, 2 for
BIN, 3 for MID, and 5 for directory. \\
\novacode{FIO\_NAME}, \novacode{FIO\_NAMELEN} & Display name for the current
entry. Some file types append useful load metadata to the display string. \\
\novacode{FIO\_SIZEL/H}, \novacode{FIO\_SIZE2} & 24-bit data size for the
entry. BASIC and BIN sizes exclude the two-byte load-address prefix. \\
\bottomrule
\end{tabularx}
\end{center}

\novacode{fio\_pwd}, \novacode{fio\_cd}, \novacode{fio\_mkdir}, and
\novacode{fio\_rmdir} use the same filename buffer. \novacode{fio\_delete}
routes through the mounted disk file delete path. File commands operate on the
current mounted NDI device, or on an explicit prefix such as \novacode{FD1:} or
\novacode{HD0:}. If no disk is mounted, the host reports
\novacode{FIO\_ERR\_NOTMOUNTED}; if a parent directory is missing, it reports
\novacode{FIO\_ERR\_NOTFOUND}. \novacode{fio\_mkdir} and
\novacode{fio\_rmdir} are the directory-creation/removal operations; saves do
not create missing parent directories implicitly.

\section{Runtime Services}

\novacode{fio\_load\_runtime} loads a 16K runtime ROM image into the primary
\novacode{\$C000} ROM bank. Code issuing this command must be running from RAM,
because the host can replace the ROM bank while the command is in progress.

\novacode{fio\_rng} asks NovaHost for 32 random bits and returns them in
\novacode{FIO\_RNG0} through \novacode{FIO\_RNG3}. These aliases overlap
\novacode{FIO\_SRCL/H} and \novacode{FIO\_ENDL/H}, so consume the random value
before preparing another command.

\section{Using FIO Well}

\begin{itemize}
\item Always branch on the A result from the helper, then read
\novacode{FIO\_ERRCODE} before issuing another command.
\item Use \novacode{fio\_copy\_name} for pointer-based filenames; it catches the
empty and too-long cases before the host sees them.
\item Use the bulk paths for assets. Load graphics into VGC memory with
\novacode{fio\_gload}, stage large data in XRAM with XLOAD, and use the pager
when only one slice of a larger file is needed.
\item Do not assume \novacode{FIO\_NAME} survives across subsystems. Directory
reads, current-directory queries, and other host services all reuse the same
buffer.
\item Clear stale host-visible errors with \novacode{fio\_clear\_error} before
returning to code that polls FIO status independently.
\end{itemize}
